\metatitle{Independence and circuits in colored matroids}
\metadescription{Partial colorings, circuits, and the recovery of ordinary
matroid independence and duality in two colors.}

\section[independenceAndCircuits]{Independence and circuits}

For an ordinary matroid, every independent set extends to a basis, and the
circuits are the minimal sets that do not. The same definitions apply to
partial colorings, where we ask whether the prescribed colors extend to a
word of the colored matroid.

Let $\mathcal F$ be a colored matroid on $E$ with colors $1,\dotsc,k$. For a
word $u\in\mathcal F$, write
\[
 \defin{\widehat B(u)}=\{(e,u_e):e\in E\}.
\]

\begin{definition}[Independent partial assignments and circuits]
A \defin{partial assignment} is a set
$\tau\subseteq E\times[k]$ containing at most one pair with first coordinate
$e$, for every $e\in E$. It is \defin{independent} if
\[
 \tau\subseteq\widehat B(u)
\]
for some $u\in\mathcal F$. A \defin{circuit} is an inclusion-minimal partial
assignment that is not independent.
\end{definition}

Every subset of an independent partial assignment is independent. The
maximal independent partial assignments are exactly the sets
$\widehat B(u)$, for $u\in\mathcal F$.

\subsection[twoColorIndependence]{Two colors}

The preceding definitions recover ordinary matroid independence without any
change of meaning. Suppose that the colors are $1$ and $0$, and let $M$ be
the matroid whose bases are the color-$1$ classes. Its dual $M^*$ has the
color-$0$ classes as bases.

\begin{proposition}[Classical matroid independence]
Let $\tau$ prescribe color $1$ on $I$ and color $0$ on $J$, where
$I,J\subseteq E$ are disjoint. Then $\tau$ is independent if and only if
\[
 I\text{ is independent in }M
 \qquad\text{and}\qquad
 J\text{ is independent in }M^*.
\]
Consequently, every circuit uses a single color. The color-$1$ circuits are
the circuits of $M$, while the color-$0$ circuits are the circuits of $M^*$,
or equivalently the cocircuits of $M$.
\end{proposition}

\begin{proof}
If $\tau$ extends to a word, then $I$ lies in its color-$1$ basis and $J$
lies in its complementary color-$0$ basis. Thus the two sets are independent
in $M$ and $M^*$, respectively.

Conversely, suppose that $I$ is independent in $M$ and $J$ is independent
in $M^*$. Some basis of $M^*$ contains $J$, so its complementary basis of
$M$ avoids $J$. Hence the restriction of $M$ to $E\setminus J$ has full
rank. Extend $I$ to a basis of this restriction. Its indicator word extends
$\tau$.

It follows that a nonextendible partial assignment contains either a circuit
of $M$ colored $1$ or a circuit of $M^*$ colored $0$. Minimality gives the
last assertion.
\end{proof}

Taking $J=\varnothing$ shows that a set of positions colored $1$ is
independent precisely when it is an ordinary independent set of $M$. Taking
$I=\varnothing$ gives the corresponding statement for the dual. Thus bases,
independent sets, circuits, cocircuits, and duality all occur inside the
two-color theory.

\subsection[mixedCircuits]{Mixed circuits}

For three or more colors, the individual color projections do not contain
all the independence information.

\begin{proposition}[Monochromatic circuits]
Fix a color $c$. A circuit whose assignments all have color $c$ is exactly
an ordinary circuit of the color-$c$ projection matroid, with all its
positions assigned color $c$.
\end{proposition}

\begin{proof}
A set $I$ is independent in the color-$c$ projection if and only if it is
contained in $B(u,c)$ for some $u\in\mathcal F$. This is equivalent to the
partial assignment $I\times\{c\}$ being independent. Minimal nonextendible
sets therefore agree on both sides.
\end{proof}

Mixed circuits record compatibility between several color classes. Consider
the seven-word colored matroid
\[
 \mathcal F={\mathtt{1123},\mathtt{1213},\mathtt{1321},
 \mathtt{2113},\mathtt{3121},\mathtt{2311},\mathtt{3211}\}.
\]

\begin{example}[A mixed circuit]
The partial assignment
\[
 \{(1,3),(2,1),(3,1)\}
\]
is a circuit. The two occurrences of color $1$ force the remaining colors
$2$ and $3$ onto positions $1$ and $4$, but color $3$ is already prescribed
at position $1$ and no word places color $2$ at position $4$.

Removing the assignments at positions $1$, $2$, and $3$ gives partial
assignments extending to $\mathtt{2113}$, $\mathtt{3211}$, and
$\mathtt{3121}$, respectively. The color-$1$ positions $\{2,3\}$ and the
color-$3$ position $\{1\}$ are separately independent in their projection
matroids; their joint prescription is not.
\end{example}

\subsection[conditionalCircuits]{Conditional circuits}

Mixed circuits become ordinary circuits after the other colors have been
fixed.

\begin{proposition}[Circuits after pinning]
Let $C$ be a circuit, and let $c$ be a color occurring in $C$. Put
\[
 I=\{e:(e,c)\in C\},
 \qquad
 \tau=C\setminus(I\times\{c\}).
\]
Pin the assignments in $\tau$ and take the color-$c$ projection of the
remaining family. Then $I$ is a circuit of this ordinary matroid.
\end{proposition}

\begin{proof}
The proper subset $\tau$ extends to a word. No completion after pinning
$\tau$ can contain every position of $I$ in color $c$, since that would
extend $C$. If $e\in I$, minimality of $C$ gives a completion of
$C\setminus\{(e,c)\}$. Thus $I\setminus\{e\}$ is independent after pinning
$\tau$. Hence $I$ is an ordinary circuit in the conditional color
projection.
\end{proof}

There is also a direct remnant of basis exchange for arbitrary partial
assignments.

\begin{proposition}[Adding one assignment]
Suppose $\sigma\subseteq\tau$ are independent partial assignments and
$a\notin\tau$, while $\sigma\cup\{a\}$ is independent. We can remove at most
one assignment of each color from $\tau$ so that $a$ can be added and the
result remains independent.
\end{proposition}

\begin{proof}
Choose a word $u$ extending $\tau$ and a word $v$ extending
$\sigma\cup\{a\}$. If $u$ already contains $a$, there is nothing to prove.
Otherwise, apply symmetric cycle exchange at the position of $a$. The two
words agree on every assignment of $\sigma$, so the exchange avoids those
positions. Remove from $\tau$ the assignments changed in the first word.
The resulting partial assignment together with $a$ extends to the exchanged
word. Since the exchange cycle uses each color at most once, at most one
assignment of each color was removed.
\end{proof}
